\chapter{Carotid Angioplasty and Stenting}

\section*{What Is This Test or Procedure?}
Carotid angioplasty and stenting (CAS) is a minimally invasive procedure that opens a narrowed carotid artery using a balloon and stent, usually placed through a transfemoral approach.

\section*{Alternative Names}
Carotid artery stenting; CAS; percutaneous carotid revascularization

\section*{Principle}
A balloon catheter is advanced into the carotid stenosis, inflated to dilate the vessel, and a self-expanding stent is deployed to scaffold the artery open. An embolic protection device is often used to capture debris released during dilation.

\section*{History}
Carotid stenting was developed in the 1990s as an alternative to endarterectomy. It proved particularly useful for patients at high surgical risk, and randomized trials established its role in the 2000s.

\section*{Why This Test or Procedure Is Ordered}
Performed for symptomatic carotid stenosis (greater than 50 percent) or asymptomatic stenosis (greater than 60-80 percent) in patients at high risk for open surgery, or for restenosis after endarterectomy.

\section*{Is This a Primary or Secondary Test or Procedure}
Primary or alternative revascularization for carotid stenosis; secondary after failed endarterectomy.

\section*{How This Test or Procedure Is Performed}
Performed under local or general anesthesia in an angiography suite. A sheath is placed in the femoral artery, a guidewire crosses the stenosis, an embolic protection device is deployed, the balloon is inflated, and the stent is placed. Angiography confirms patency.

\section*{What Preparation Is Needed}
Preoperative carotid ultrasound and often CT or MR angiography, dual antiplatelet therapy (aspirin plus clopidogrel) started beforehand, and assessment of renal function and contrast allergy.

\section*{Contraindications and Precautions}
Severe aortic or carotid tortuosity, heavily calcified or thrombotic lesions, or recent large stroke are relative contraindications. Advanced age and unfavorable arch anatomy increase stroke risk.

\section*{Risks}
Periprocedural stroke or transient ischemic attack, dissection, stent thrombosis or restenosis, bleeding or pseudoaneurysm at the puncture site, and contrast nephropathy.

\section*{What Happens After the Test or Procedure}
Observation with neurologic checks, dual antiplatelet therapy continued for weeks to months, and puncture-site care. Most patients are discharged within 1-2 days.

\section*{Understanding Your Results}
Post-procedure angiography and ultrasound confirm stent patency and residual stenosis. Results guide antiplatelet duration and risk-factor management.

\section*{Normal Results}
A patent stent with less than 30 percent residual stenosis, no new neurologic deficit, and no access-site complication.

\section*{Follow-up Tests and Procedures}
Carotid ultrasound at 1, 6, and 12 months and then annually to detect in-stent restenosis, along with aggressive risk-factor control.

\section*{References}
\begin{enumerate}
  \item MedlinePlus. Carotid artery stenting. U.S. National Library of Medicine; 2025.
  \item Current Medical Diagnosis and Treatment. McGraw-Hill; 2025.
\end{enumerate}

\clearpage
